% SPDX-License-Identifier: Apache-2.0
% covers: Node
\datasheet{Node}{A node of the network on chip: a switch per virtual channel.}

\dsfacts{\code{lib/parts/src/bus/noc/node.rs}}{\code{txhdl\_parts::bus::noc::node::Node}}%
{\code{//docs:noc}}

\subsection*{Function}

A node is two \code{Switch} units: \code{req}, for what a host sends
and a peripheral receives, and \code{rsp}, for what a peripheral sends
and a host receives. The two share nothing. So a request held up behind
a full buffer cannot hold up the response that would empty it.

The node's place in the lattice is two inputs, \code{col} and
\code{row}, which it passes to both switches. Whoever places the node
holds them at constants with \code{tie}, so every node of a lattice is
one type (issue 635), and \code{//docs:noc} states that synthesis
folds the constants into the comparisons that route a packet.

Its other twenty ports are written out one by one. A port name starts
with \code{q} for the request channel or \code{p} for the response
channel. The side follows: \code{n}, \code{s}, \code{w}, \code{e}, or
\code{x} for the exit. Last is \code{i} or \code{o}. So \code{qei} is
the request input from the east, and \code{pxo} is the response output
to the exit.

\subsection*{Parameters}

\begin{center}\footnotesize
\begin{tabular}{@{}lp{0.7\textwidth}@{}}
\toprule
Parameter & Meaning \\
\midrule
\code{XB} & The width of a column coordinate, and of \code{col}, in
bits. A lattice is at most \code{1 << XB} columns wide. \\
\code{YB} & The width of a row coordinate, and of \code{row}, in bits.
A lattice is at most \code{1 << YB} rows high. \\
\code{A} & The address width of the AXI link. \\
\code{D} & The data width of the AXI link. \\
\code{S} & The strobe width, which is \code{D / 8}. \\
\code{I} & The identifier width of the AXI link. \\
\bottomrule
\end{tabular}
\end{center}

The tables below are for \code{Node<2, 2, 32, 32, 4, 2>}, the node on
the link that \code{//soc} uses.

\dstables{Node}

\subsection*{Behaviour}

\begin{itemize}
\item The node has no logic of its own. It runs \code{req} on the ten
\code{q} ports and \code{rsp} on the ten \code{p} ports, side by side,
both at the place \code{col} and \code{row} say. What it holds is its
two switches' state, which is six bits: a \code{turn} of three in
each.
\item Each switch routes in dimension order and moves one packet per
cycle, choosing among the inputs that can move by a round robin of its
own, as the \code{Switch} sheet states. The two switches take turns
independently: a busy request channel does not decide which response
moves.
\item A host's bridge sends requests into \code{qxi} and takes
responses from \code{pxo}. A peripheral's bridge takes requests from
\code{qxo} and sends responses into \code{pxi}.
\end{itemize}

\subsection*{Verification}

\begin{itemize}
\item Eight tests in \code{lib/parts/src/bus/noc.rs}, in
\code{//lib/parts:parts\_test}, run four nodes on a two by two lattice
with both bridges, both AXI trackers and memory. Among them,
\code{a\_burst\_crosses\_the\_lattice\_and\_its\_answer\_comes\_back}
writes a word from the host at 0,0 to a memory at 1,1 and reads it
back; \code{two\_hosts\_share\_a\_memory\_and\_each\_answer\_goes\_home}
puts hosts at 0,0 and 1,0 on that memory; and
\code{four\_peripherals\_share\_one\_node} puts a router of four
memories behind the node at 1,1. The others cross the lattice with
reads of four beats, writes of two and of sixteen, a fixed write and a
wrapping one.
\item \code{a\_core\_draws\_a\_solid\_and\_a\_gpu\_fills\_it}, in
\code{//soc:demo\_test}, runs the core, the GPU, a memory and a serial
port on the four corners of one lattice.
\item The node's netlist is checked inside \code{Mesh}: \code{ex\_mesh}
lowers a mesh of six nodes, and \code{//docs:mesh\_sim\_mesh\_tb\_test}
and \code{//docs:mesh\_vsim\_test} run it under nvc and Verilator
against the Rust run, at the mesh's ports.
\end{itemize}

\subsection*{Used by}

The tests of the network in \code{lib/parts/src/bus/noc.rs}, and
\code{soc::run} in \code{soc/src/lib.rs}, each with four nodes wired by
\code{mesh::lattice}; and \code{Mesh}, which holds \code{N} of them.

\subsection*{Limits}

\code{mesh::lattice} makes the channels in simulation and does not
lower; \code{Mesh} is the same wiring as a unit that does. On a lattice
that \code{mesh::lattice} makes, a link at the edge drives a channel
nobody reads. Each of the two switches moves one packet per cycle, not
one per port.
